\documentclass{beamer}

\usepackage{beamerthemeSingapore}
\usepackage[brazil]{babel}
\usepackage[T1]{fontenc}
\usepackage[utf8]{inputenc}
\usepackage{graphicx}
\usepackage{url}
\usepackage{xcolor}
\usepackage{booktabs}
\usepackage{array}
\usepackage{ragged2e}
\usepackage{etoolbox}
\usepackage{amsmath}
\usepackage{minted}
\usepackage{tikz}
\usetikzlibrary{positioning, arrows.meta, calc}
% Grafo G e estilos da aula de representação e percursos, e as figuras desta
% aula, compartilhadas com o roteiro.
% Grafos da aula de representação e percursos. Compartilhado entre o roteiro
% (figuras.tex) e os slides (grafos.tex da pasta 00_grafos). Os grafos são os
% mesmos dos programas da pasta codes.

\tikzset{
  gv/.style = {circle, draw=black, fill=white, inner sep=0pt,
    minimum size=1.6em, line width=0.6pt, font=\sffamily\small},
  ga/.style = {draw=black!55, line width=0.7pt},
  garv/.style = {draw=blue!70!black, line width=2.2pt},
  gnota/.style = {font=\sffamily\scriptsize, text=blue!70!black},
  arco/.style = {-{Stealth[length=2.4mm]}, draw=black!70, line width=0.7pt},
  pv/.style = {rectangle, rounded corners=2pt, draw=black, fill=white,
    inner sep=3pt, line width=0.6pt, font=\ttfamily\small},
  arcov/.style = {-{Stealth[length=2.4mm]}, draw=red!75!black, line width=1.6pt}
}

% ------------------------------------------------------------------
% Grafo G, de bfs.go, dfs.go e representacao.go.
\newcommand{\Gnos}{%
  \node [gv] (g0) at (0, 1.5) {0};
  \node [gv] (g1) at (1.5, 3) {1};
  \node [gv] (g2) at (1.5, 1.5) {2};
  \node [gv] (g5) at (1.5, 0) {5};
  \node [gv] (g3) at (3, 3) {3};
  \node [gv] (g4) at (3, 0) {4};
  \node [gv] (g6) at (4.5, 1.5) {6};
  \node [gv] (g7) at (6, 1.5) {7};
}
\newcommand{\Garestas}{%
  \foreach \a/\b in {0/1, 0/2, 0/5, 1/3, 2/3, 2/4, 3/6, 4/5, 4/6, 6/7}
    \draw [ga] (g\a) -- (g\b);
}

\newcommand{\grafoG}{%
\begin{tikzpicture}
  \Gnos
  \Garestas
\end{tikzpicture}}

% Árvore da busca em largura a partir de 0, com a distância de cada vértice.
\newcommand{\grafoGbfs}{%
\begin{tikzpicture}
  \Gnos
  \Garestas
  \foreach \a/\b in {0/1, 0/2, 0/5, 1/3, 2/4, 3/6, 6/7}
    \draw [garv] (g\a) -- (g\b);
  \Gnos
  \foreach \v/\d/\p in {0/0/left, 1/1/above, 2/1/right, 5/1/below,
                        3/2/above, 4/2/below, 6/3/below, 7/4/above}
    \node [gnota, \p=1pt of g\v] {d=\d};
\end{tikzpicture}}

% Árvore da busca em profundidade a partir de 0, com a ordem de visita.
\newcommand{\grafoGdfs}{%
\begin{tikzpicture}
  \Gnos
  \Garestas
  \foreach \a/\b in {0/1, 1/3, 3/2, 2/4, 4/5, 4/6, 6/7}
    \draw [garv] (g\a) -- (g\b);
  \Gnos
  \foreach \v/\d/\p in {0/1/left, 1/2/above, 3/3/above, 2/4/right,
                        4/5/below, 5/6/below, 6/7/below, 7/8/above}
    \node [gnota, \p=1pt of g\v] {\d\textordmasculine};
\end{tikzpicture}}

% ------------------------------------------------------------------
% Grafo de componentes.go: 10 vértices, 4 componentes.
\newcommand{\grafoComp}{%
\begin{tikzpicture}
  \node [gv] (h0) at (0, 1.6) {0};
  \node [gv] (h1) at (-0.7, 0.4) {1};
  \node [gv] (h2) at (0.7, 0.4) {2};
  \node [gv] (h3) at (2.2, 1.6) {3};
  \node [gv] (h4) at (2.2, 0.4) {4};
  \node [gv] (h5) at (3.7, 1.6) {5};
  \node [gv] (h6) at (5.0, 1.6) {6};
  \node [gv] (h7) at (3.7, 0.4) {7};
  \node [gv] (h8) at (5.0, 0.4) {8};
  \node [gv] (h9) at (6.5, 1.0) {9};
  \foreach \a/\b in {0/1, 0/2, 1/2, 3/4, 5/6, 5/7, 6/8, 7/8}
    \draw [ga] (h\a) -- (h\b);
  \node [gnota, below=0.25cm of h1, xshift=0.7cm] {id 0};
  \node [gnota, below=0.25cm of h4] {id 1};
  \node [gnota, below=0.25cm of h7, xshift=0.65cm] {id 2};
  \node [gnota, below=0.85cm of h9] {id 3};
\end{tikzpicture}}

% ------------------------------------------------------------------
% Pacotes de topologica.go. O arco B -> A significa que A importa B.
\newcommand{\Pnos}{%
  \node [pv] (modelo) at (0, 2.4) {modelo};
  \node [pv] (log) at (0, 1.2) {log};
  \node [pv] (config) at (0, 0) {config};
  \node [pv] (banco) at (2.6, 1.2) {banco};
  \node [pv] (servico) at (5.2, 2.4) {servico};
  \node [pv] (http) at (7.8, 1.2) {http};
  \node [pv] (main) at (10.4, 1.2) {main};
}
\newcommand{\Parcos}{%
  \draw [arco] (modelo) -- (banco);
  \draw [arco] (log) -- (banco);
  \draw [arco] (config) -- (banco);
  \draw [arco] (banco) -- (servico);
  \draw [arco] (modelo) -- (servico);
  \draw [arco] (log) -- (servico);
  \draw [arco] (log) to [bend right=14] (http);
  \draw [arco] (servico) -- (http);
  \draw [arco] (servico) -- (main);
  \draw [arco] (http) -- (main);
  \draw [arco] (config) to [bend right=8] (main);
}
\newcommand{\grafoPacotes}{%
\begin{tikzpicture}
  \Pnos
  \Parcos
\end{tikzpicture}}

% O mesmo grafo com o arco servico -> modelo, que fecha um ciclo.
\newcommand{\grafoPacotesCiclo}{%
\begin{tikzpicture}
  \Pnos
  \Parcos
  \draw [arcov] (modelo) -- (banco);
  \draw [arcov] (banco) -- (servico);
  \draw [arcov] (servico) to [bend right=25] (modelo);
\end{tikzpicture}}

% A ordem topológica encontrada, da esquerda para a direita: todo arco aponta
% para a direita.
\newcommand{\grafoPacotesOrdem}{%
\begin{tikzpicture}
  \foreach \n/\x in {log/0, config/1.7, modelo/3.4, banco/5.1,
                     servico/6.8, http/8.5, main/10.2}
    \node [pv] (\n) at (\x, 0) {\n};
  \draw [arco] (modelo) -- (banco);
  \draw [arco] (banco) -- (servico);
  \draw [arco] (servico) -- (http);
  \draw [arco] (http) -- (main);
  \draw [arco] (log.north) .. controls +(0, 1.0) and +(0, 1.0) .. (banco.north);
  \draw [arco] (log.north) .. controls +(0, 1.4) and +(0, 1.4) .. (servico.north);
  \draw [arco] (log.north) .. controls +(0, 1.8) and +(0, 1.8) .. (http.north);
  \draw [arco] (servico.north) .. controls +(0, 0.8) and +(0, 0.8) .. (main.north);
  \draw [arco] (config.south) .. controls +(0, -0.8) and +(0, -0.8) .. (banco.south);
  \draw [arco] (modelo.south) .. controls +(0, -1.1) and +(0, -1.1) .. (servico.south);
  \draw [arco] (config.south) .. controls +(0, -1.6) and +(0, -1.6) .. (main.south);
\end{tikzpicture}}

% ------------------------------------------------------------------
% Componentes fortemente conexas: três componentes, destacadas pelo fundo.
\newcommand{\grafoCFC}{%
\begin{tikzpicture}
  \fill [orange!18, rounded corners=8pt] (-0.55, -0.6) rectangle (2.05, 2.1);
  \fill [green!18, rounded corners=8pt] (3.15, -0.6) rectangle (5.75, 2.1);
  \fill [blue!12, rounded corners=8pt] (6.85, 0.15) rectangle (7.95, 1.35);
  \node [gv] (c0) at (0, 1.5) {0};
  \node [gv] (c1) at (1.5, 1.5) {1};
  \node [gv] (c2) at (0.75, 0) {2};
  \node [gv] (c3) at (3.7, 1.5) {3};
  \node [gv] (c4) at (5.2, 1.5) {4};
  \node [gv] (c5) at (4.45, 0) {5};
  \node [gv] (c6) at (7.4, 0.75) {6};
  \draw [arco] (c0) -- (c1);
  \draw [arco] (c1) -- (c2);
  \draw [arco] (c2) -- (c0);
  \draw [arco] (c3) -- (c4);
  \draw [arco] (c4) -- (c5);
  \draw [arco] (c5) -- (c3);
  \draw [arco] (c1) -- (c3);
  \draw [arco] (c2) -- (c5);
  \draw [arco] (c4) -- (c6);
\end{tikzpicture}}

% ------------------------------------------------------------------
% Execução passo a passo das duas buscas em G, a partir de 0. Os estados
% seguem a saída de bfs.go e dfs.go. Cores: branco, ainda não alcançado;
% cinza, na fila (largura) ou com a chamada em andamento (profundidade); preto,
% já retirado da fila ou com a chamada terminada. A borda vermelha marca o
% vértice em processamento. Em azul, as arestas da árvore; em laranja, as
% arestas que entraram na árvore no passo.
\tikzset{
  gcinza/.style = {fill=black!25},
  gpreto/.style = {fill=black!80, text=white},
  gatual/.style = {draw=red!80!black, line width=2pt},
  gnova/.style = {draw=orange!90!black, line width=2.6pt},
  gcaixa/.style = {rectangle, draw=black, fill=black!25, line width=0.6pt,
    minimum width=0.55cm, minimum height=0.5cm, inner sep=0pt,
    font=\sffamily\small},
  gquadro/.style = {rectangle, draw=black, fill=black!25, line width=0.6pt,
    minimum width=1.7cm, minimum height=0.36cm, inner sep=0pt,
    font=\ttfamily\scriptsize},
  gtexto/.style = {anchor=west, font=\sffamily\footnotesize, inner sep=0pt},
  gtabela/.style = {anchor=east, font=\ttfamily\small, inner sep=0pt},
  glegenda/.style = {anchor=west, font=\sffamily\scriptsize, inner sep=0pt}
}

% Posição do rótulo de distância em volta de cada vértice, em graus.
\newcommand{\Gang}[1]{\ifcase#1 180\or 90\or 0\or 90\or 270\or 270\or 270\or 90\fi}

% Grafo G no estado de um passo. #1: vértice em processamento (pode ser
% vazio); #2: vértices cinza; #3: vértices pretos; #4: arestas da árvore;
% #5: arestas novas no passo.
\newcommand{\Gestado}[5]{%
  \Gnos
  \Garestas
  \foreach \a/\b in {#4} \draw [garv] (g\a) -- (g\b);
  \foreach \a/\b in {#5} \draw [gnova] (g\a) -- (g\b);
  \Gnos
  \foreach \v in {#2} \node [gv, gcinza] at (g\v) {\v};
  \foreach \v in {#3} \node [gv, gpreto] at (g\v) {\v};
  \foreach \v in {#1} \draw [gatual] (g\v) circle (0.8em);
}

% Um passo da busca em largura. #1 a #5 como em \Gestado; #6: fila, do início
% para o fim; #7: distâncias conhecidas, em pares v/d; #8: o que o passo faz.
\newcommand{\passoBFS}[8]{%
\begin{tikzpicture}
  \useasboundingbox (-1.1, -2.55) rectangle (8.8, 3.65);
  \Gestado{#1}{#2}{#3}{#4}{#5}
  \foreach \v/\d in {#7} {
    \edef\ang{\Gang{\v}}
    \path (g\v) ++(\ang:0.66 and 0.43) node [gnota] {d=\d};
  }
  \node [gtabela] at (7.0, 3.4) {v};
  \node [gtabela] at (7.95, 3.4) {dist};
  \node [gtabela] at (8.75, 3.4) {pai};
  \draw [black!60] (6.6, 3.18) -- (8.8, 3.18);
  \foreach \v in {0, ..., 7}
    \node [gtabela] at (7.0, 2.92 - 0.38*\v) {\v};
  \foreach \v/\d in {#7}
    \node [gtabela] at (7.95, 2.92 - 0.38*\v) {\d};
  \node [gtabela] at (8.75, 2.92) {-1};
  \foreach \a/\b in {#4} \node [gtabela] at (8.75, 2.92 - 0.38*\b) {\a};
  \foreach \a/\b in {#5} \node [gtabela] at (8.75, 2.92 - 0.38*\b) {\a};
  \node [gtexto] at (-1.1, -1.05) {fila:};
  \foreach \v [count=\i] in {#6}
    \node [gcaixa] at (-0.2 + 0.6*\i, -1.05) {\v};
  \node [gtexto] at (-1.1, -1.75) {#8};
  \node [gv, gcinza, minimum size=1em] at (-0.9, -2.35) {};
  \node [glegenda] at (-0.65, -2.35) {na fila};
  \node [gv, gpreto, minimum size=1em] at (0.9, -2.35) {};
  \node [glegenda] at (1.15, -2.35) {já retirado};
  \node [gv, gatual, minimum size=1em] at (3.2, -2.35) {};
  \node [glegenda] at (3.45, -2.35) {retirado agora};
\end{tikzpicture}}

% Um passo da busca em profundidade. #1 a #5 como em \Gestado, com #2 também
% sendo a pilha de chamadas, da base para o topo; #6: pré-ordem até o passo;
% #7: pós-ordem até o passo; #8: o que o passo faz.
\newcommand{\passoDFS}[8]{%
\begin{tikzpicture}
  \useasboundingbox (-1.1, -2.55) rectangle (8.8, 3.65);
  \Gestado{#1}{#2}{#3}{#4}{#5}
  \gdef\Gtopo{0}
  \foreach \v [count=\i] in {#2} {
    \node [gquadro] at (7.9, -0.2 + 0.42*\i) {visita(\v)};
    \xdef\Gtopo{\i}
  }
  \ifnum\Gtopo>0
    \node [gquadro, gatual, fill=none] at (7.9, -0.2 + 0.42*\Gtopo) {};
  \fi
  \node [glegenda, anchor=center] at (7.9, -0.25) {pilha de chamadas};
  \node [gtexto] at (-1.1, -0.85) {pré-ordem: \texttt{#6}};
  \node [gtexto] at (-1.1, -1.3) {pós-ordem: \texttt{#7}};
  \node [gtexto] at (-1.1, -1.75) {#8};
  \node [gv, gcinza, minimum size=1em] at (-0.9, -2.35) {};
  \node [glegenda] at (-0.65, -2.35) {chamada em andamento};
  \node [gv, gpreto, minimum size=1em] at (3.0, -2.35) {};
  \node [glegenda] at (3.25, -2.35) {chamada terminada};
  \node [gv, gatual, minimum size=1em] at (6.4, -2.35) {};
  \node [glegenda] at (6.65, -2.35) {em execução};
\end{tikzpicture}}

% Passos numerados a partir de 1, para \only nos slides.
\newcommand{\bfsPasso}[1]{\ifcase#1
\or\passoBFS{}{0}{}{}{}{0}{0/0}{início: marca o 0 e o põe na fila}
\or\passoBFS{0}{1, 2, 5}{}{}{0/1, 0/2, 0/5}{1, 2, 5}{0/0, 1/1, 2/1, 5/1}{retira 0, lista \texttt{[1 2 5]}: 1, 2 e 5 entram na fila}
\or\passoBFS{1}{2, 5, 3}{0}{0/1, 0/2, 0/5}{1/3}{2, 5, 3}{0/0, 1/1, 2/1, 3/2, 5/1}{retira 1, lista \texttt{[0 3]}: 3 entra na fila}
\or\passoBFS{2}{5, 3, 4}{0, 1}{0/1, 0/2, 0/5, 1/3}{2/4}{5, 3, 4}{0/0, 1/1, 2/1, 3/2, 4/2, 5/1}{retira 2, lista \texttt{[0 3 4]}: 4 entra na fila}
\or\passoBFS{5}{3, 4}{0, 1, 2}{0/1, 0/2, 0/5, 1/3, 2/4}{}{3, 4}{0/0, 1/1, 2/1, 3/2, 4/2, 5/1}{retira 5, lista \texttt{[0 4]}: todos marcados, nada entra}
\or\passoBFS{3}{4, 6}{0, 1, 2, 5}{0/1, 0/2, 0/5, 1/3, 2/4}{3/6}{4, 6}{0/0, 1/1, 2/1, 3/2, 4/2, 5/1, 6/3}{retira 3, lista \texttt{[1 2 6]}: 6 entra na fila}
\or\passoBFS{4}{6}{0, 1, 2, 5, 3}{0/1, 0/2, 0/5, 1/3, 2/4, 3/6}{}{6}{0/0, 1/1, 2/1, 3/2, 4/2, 5/1, 6/3}{retira 4, lista \texttt{[2 5 6]}: todos marcados, nada entra}
\or\passoBFS{6}{7}{0, 1, 2, 5, 3, 4}{0/1, 0/2, 0/5, 1/3, 2/4, 3/6}{6/7}{7}{0/0, 1/1, 2/1, 3/2, 4/2, 5/1, 6/3, 7/4}{retira 6, lista \texttt{[3 4 7]}: 7 entra na fila}
\or\passoBFS{7}{}{0, 1, 2, 5, 3, 4, 6}{0/1, 0/2, 0/5, 1/3, 2/4, 3/6, 6/7}{}{}{0/0, 1/1, 2/1, 3/2, 4/2, 5/1, 6/3, 7/4}{retira 7, lista \texttt{[6]}: fila vazia, fim da busca}
\fi}
\newcommand{\dfsPasso}[1]{\ifcase#1
\or\passoDFS{0}{0}{}{}{}{0}{}{\texttt{visita(0)}: lista \texttt{[1 2 5]}, desce para o 1}
\or\passoDFS{1}{0, 1}{}{}{0/1}{0 1}{}{\texttt{visita(1)}: lista \texttt{[0 3]}, 0 já marcado, desce para o 3}
\or\passoDFS{3}{0, 1, 3}{}{0/1}{1/3}{0 1 3}{}{\texttt{visita(3)}: lista \texttt{[1 2 6]}, 1 já marcado, desce para o 2}
\or\passoDFS{2}{0, 1, 3, 2}{}{0/1, 1/3}{3/2}{0 1 3 2}{}{\texttt{visita(2)}: lista \texttt{[0 3 4]}, 0 e 3 já marcados, desce para o 4}
\or\passoDFS{4}{0, 1, 3, 2, 4}{}{0/1, 1/3, 3/2}{2/4}{0 1 3 2 4}{}{\texttt{visita(4)}: lista \texttt{[2 5 6]}, 2 já marcado, desce para o 5}
\or\passoDFS{5}{0, 1, 3, 2, 4, 5}{}{0/1, 1/3, 3/2, 2/4}{4/5}{0 1 3 2 4 5}{}{\texttt{visita(5)}: lista \texttt{[0 4]}, todos marcados}
\or\passoDFS{4}{0, 1, 3, 2, 4}{5}{0/1, 1/3, 3/2, 2/4, 4/5}{}{0 1 3 2 4 5}{5}{\texttt{termina 5}: volta ao 4; o 6 não está marcado}
\or\passoDFS{6}{0, 1, 3, 2, 4, 6}{5}{0/1, 1/3, 3/2, 2/4, 4/5}{4/6}{0 1 3 2 4 5 6}{5}{\texttt{visita(6)}: lista \texttt{[3 4 7]}, 3 e 4 já marcados, desce para o 7}
\or\passoDFS{7}{0, 1, 3, 2, 4, 6, 7}{5}{0/1, 1/3, 3/2, 2/4, 4/5, 4/6}{6/7}{0 1 3 2 4 5 6 7}{5}{\texttt{visita(7)}: lista \texttt{[6]}, todos marcados}
\or\passoDFS{6}{0, 1, 3, 2, 4, 6}{5, 7}{0/1, 1/3, 3/2, 2/4, 4/5, 4/6, 6/7}{}{0 1 3 2 4 5 6 7}{5 7}{\texttt{termina 7}: volta ao 6; fim da lista}
\or\passoDFS{4}{0, 1, 3, 2, 4}{5, 7, 6}{0/1, 1/3, 3/2, 2/4, 4/5, 4/6, 6/7}{}{0 1 3 2 4 5 6 7}{5 7 6}{\texttt{termina 6}: volta ao 4; fim da lista}
\or\passoDFS{2}{0, 1, 3, 2}{5, 7, 6, 4}{0/1, 1/3, 3/2, 2/4, 4/5, 4/6, 6/7}{}{0 1 3 2 4 5 6 7}{5 7 6 4}{\texttt{termina 4}: volta ao 2; fim da lista}
\or\passoDFS{3}{0, 1, 3}{5, 7, 6, 4, 2}{0/1, 1/3, 3/2, 2/4, 4/5, 4/6, 6/7}{}{0 1 3 2 4 5 6 7}{5 7 6 4 2}{\texttt{termina 2}: volta ao 3; 6 já marcado, fim da lista}
\or\passoDFS{1}{0, 1}{5, 7, 6, 4, 2, 3}{0/1, 1/3, 3/2, 2/4, 4/5, 4/6, 6/7}{}{0 1 3 2 4 5 6 7}{5 7 6 4 2 3}{\texttt{termina 3}: volta ao 1; fim da lista}
\or\passoDFS{0}{0}{5, 7, 6, 4, 2, 3, 1}{0/1, 1/3, 3/2, 2/4, 4/5, 4/6, 6/7}{}{0 1 3 2 4 5 6 7}{5 7 6 4 2 3 1}{\texttt{termina 1}: volta ao 0; 2 e 5 já marcados, fim da lista}
\or\passoDFS{}{}{5, 7, 6, 4, 2, 3, 1, 0}{0/1, 1/3, 3/2, 2/4, 4/5, 4/6, 6/7}{}{0 1 3 2 4 5 6 7}{5 7 6 4 2 3 1 0}{\texttt{termina 0}: fim da busca}
\fi}

% ------------------------------------------------------------------
% Balde de tinta: imagem de 10 por 7 pixels, com 4 vizinhos por pixel. #1 = 0
% mostra a imagem antes do clique; #1 = 1, depois, com a componente do pixel
% clicado pintada.
\newcommand{\imagemBalde}[1]{%
\begin{tikzpicture}[x=0.38cm, y=0.38cm]
  \foreach \x/\y in {2/2, 3/2, 4/2, 5/2, 6/2, 7/2, 2/3, 3/3, 6/3, 7/3, 2/4, 3/4, 4/4, 5/4, 6/4, 7/4}
    {\ifnum#1=1 \fill [orange!70] (\x, -\y) rectangle ++(1, -1);\fi}
  \foreach \x/\y in {2/1, 3/1, 4/1, 5/1, 6/1, 7/1, 1/2, 8/2, 1/3, 4/3, 5/3, 8/3, 1/4, 8/4, 2/5, 3/5, 4/5, 5/5, 6/5, 7/5}
    \fill [black!80] (\x, -\y) rectangle ++(1, -1);
  \draw [black!25, line width=0.3pt] (0, 0) grid (10, -7);
  \draw [black, line width=0.6pt] (0, 0) rectangle (10, -7);
  \fill [red!80!black] (2.5, -3.5) circle (0.18);
  \node [font=\sffamily\scriptsize, anchor=north] at (5, -7.2)
    {\ifnum#1=1 depois do clique\else antes do clique\fi};
\end{tikzpicture}}

% ------------------------------------------------------------------
% Memória de um programa: objetos e ponteiros. Os objetos que as raízes não
% alcançam são lixo.
\newcommand{\grafoMemoria}{%
\begin{tikzpicture}
  \node [pv, font=\sffamily\small] (raiz) at (0, 1.2) {raízes};
  \node [gv] (a) at (2, 2) {a};
  \node [gv] (b) at (2, 0.4) {b};
  \node [gv] (c) at (3.6, 2) {c};
  \node [gv] (d) at (3.6, 0.4) {d};
  \node [gv, gcinza, dashed] (e) at (5.6, 2) {e};
  \node [gv, gcinza, dashed] (f) at (5.6, 0.4) {f};
  \node [gv, gcinza, dashed] (g) at (7.0, 1.2) {g};
  \draw [arco] (raiz) -- (a);
  \draw [arco] (raiz) -- (b);
  \draw [arco] (a) -- (c);
  \draw [arco] (b) -- (d);
  \draw [arco] (d) -- (c);
  \draw [arco] (e) -- (f);
  \draw [arco] (f) -- (g);
  \draw [arco] (e) -- (c);
  \node [gnota, text=black!60, above=0.1cm of e, xshift=0.7cm] {lixo};
\end{tikzpicture}}

% Figuras da aula de árvores geradoras mínimas. Compartilhado entre o roteiro
% (figuras.tex) e os slides (agm.tex da pasta 00_arvores_geradoras). Usa o
% grafo G e os estilos de ../../12/img/grafos.tex, que precisa ser carregado
% antes. Os pesos são os de kruskal.go e prim.go.

\tikzset{
  gpeso/.style = {fill=white, inner sep=1pt, font=\sffamily\scriptsize},
  grej/.style = {draw=red!75!black, dashed, line width=1.6pt},
  gcorte/.style = {draw=black!60, dashed, line width=0.8pt},
  gtravessia/.style = {draw=green!50!black, line width=2.2pt},
  glinha/.style = {anchor=west, font=\ttfamily\small, inner sep=0pt}
}

% Pesos de G, no meio de cada aresta.
\newcommand{\GPpesos}{%
  \foreach \a/\b/\p in {0/1/4, 0/2/3, 0/5/7, 1/3/2, 2/3/5, 2/4/6, 3/6/9,
                        4/5/1, 4/6/8, 6/7/12}
    \path (g\a) -- node [gpeso] {\p} (g\b);
}

% Grafo G com pesos.
\newcommand{\grafoGP}{%
\begin{tikzpicture}
  \Gnos
  \Garestas
  \GPpesos
\end{tikzpicture}}

% Árvore geradora mínima de G, peso 36.
\newcommand{\grafoGPagm}{%
\begin{tikzpicture}
  \Gnos
  \Garestas
  \foreach \a/\b in {4/5, 1/3, 0/2, 0/1, 2/4, 4/6, 6/7}
    \draw [garv] (g\a) -- (g\b);
  \GPpesos
  \Gnos
\end{tikzpicture}}

% Árvore da busca em largura de G a partir de 0, com os pesos: também é uma
% árvore geradora, de peso 43.
\newcommand{\grafoGPbfs}{%
\begin{tikzpicture}
  \Gnos
  \Garestas
  \foreach \a/\b in {0/1, 0/2, 0/5, 1/3, 2/4, 3/6, 6/7}
    \draw [garv] (g\a) -- (g\b);
  \GPpesos
  \Gnos
\end{tikzpicture}}

% Um corte de G: {0, 1, 2, 3} de um lado, {4, 5, 6, 7} do outro. As arestas de
% travessia estão em verde; a mais leve delas, 2-4, está na árvore mínima.
\newcommand{\grafoGPcorte}{%
\begin{tikzpicture}
  \fill [black!10, rounded corners=8pt] (-0.5, 1.0) -- (-0.5, 3.5) -- (3.5, 3.5)
    -- (3.5, 2.4) -- (2.1, 1.0) -- cycle;
  \Gnos
  \Garestas
  \foreach \a/\b in {0/5, 2/4, 3/6} \draw [gtravessia] (g\a) -- (g\b);
  \GPpesos
  \Gnos
  \foreach \v in {0, 1, 2, 3} \node [gv, gcinza] at (g\v) {\v};
  \node [gnota, text=black!70, anchor=west] at (3.6, 3.2) {lado $S$};
  \node [gnota, text=black!70, anchor=west] at (4.6, 0.2) {lado $V - S$};
\end{tikzpicture}}

% ------------------------------------------------------------------
% Prova da propriedade do corte, no corte de \grafoGPcorte. T é uma árvore
% geradora sem e = 2-4, de peso 39. Acrescentar e fecha o ciclo
% 2-4-6-3-1-0-2, que volta ao lado S por f = 3-6. Trocar f por e dá a árvore
% mínima, de peso 36.
\tikzset{
  gciclo/.style = {draw=yellow!60, line width=8pt, line cap=round},
  grotulo/.style = {font=\sffamily\small\bfseries, inner sep=1pt}
}
% #1: arestas da árvore; #2: arestas do ciclo; #3: aresta que entra (e);
% #4: aresta que sai (f); #5: legenda.
\newcommand{\provaCorteEstado}[5]{%
  \fill [black!10, rounded corners=8pt] (-0.5, 1.0) -- (-0.5, 3.5) -- (3.5, 3.5)
    -- (3.5, 2.4) -- (2.1, 1.0) -- cycle;
  \Gnos
  \foreach \a/\b in {#2} \draw [gciclo] (g\a) -- (g\b);
  \Garestas
  \foreach \a/\b in {#1} \draw [garv] (g\a) -- (g\b);
  \foreach \a/\b in {#3} \draw [gnova] (g\a) -- (g\b);
  \foreach \a/\b in {#4} \draw [grej] (g\a) -- (g\b);
  \GPpesos
  \Gnos
  \node [grotulo, text=orange!90!black] at (2.5, 0.95) {$e$};
  \node [grotulo, text=red!75!black] at (3.95, 2.5) {$f$};
  \node [gtexto, anchor=north west, text width=6.8cm,
    execute at begin node={\hyphenpenalty=10000}] at (-0.5, -0.45) {#5};
}
\newcommand{\provaCorteN}[1]{\ifcase#1
\or\provaCorteEstado{0/1, 0/2, 1/3, 3/6, 4/5, 4/6, 6/7}{}{}{}%
  {Uma árvore geradora $T$ sem $e$ = 2-4, de peso 39. Ela atravessa o corte
   só por $f$ = 3-6.}
\or\provaCorteEstado{0/1, 0/2, 1/3, 3/6, 4/5, 4/6, 6/7}%
  {2/4, 4/6, 6/3, 3/1, 1/0, 0/2}{2/4}{}%
  {Acrescentar $e$ fecha o ciclo 2-4-6-3-1-0-2. Ele sai do lado cinza por
   $e$ e volta por $f$.}
\or\provaCorteEstado{0/1, 0/2, 1/3, 4/5, 4/6, 6/7}{}{2/4}{3/6}%
  {Retirar $f$ (peso 9) e manter $e$ (peso 6) dá outra árvore geradora, de
   peso $39 - 9 + 6 = 36$.}
\fi}
\newcommand{\provaCorte}[1]{%
\begin{tikzpicture}
  \useasboundingbox (-0.5, -1.55) rectangle (6.4, 3.55);
  \provaCorteN{#1}
\end{tikzpicture}}
% Os três passos lado a lado, para o roteiro.
\newcommand{\provaCorteTudo}{%
\begin{tikzpicture}
  \begin{scope}\provaCorteN{1}\end{scope}
  \begin{scope}[xshift=7.4cm]\provaCorteN{2}\end{scope}
  \begin{scope}[xshift=14.8cm]\provaCorteN{3}\end{scope}
\end{tikzpicture}}

% ------------------------------------------------------------------
% Mapa da atividade: a sede S e as escolas A a E, com o custo de cada ligação
% possível em km de cabo.
\newcommand{\Mnos}{%
  \node [gv] (mS) at (0, 1.5) {S};
  \node [gv] (mA) at (1.8, 3) {A};
  \node [gv] (mB) at (1.8, 0) {B};
  \node [gv] (mD) at (4.2, 3) {D};
  \node [gv] (mC) at (4.2, 0) {C};
  \node [gv] (mE) at (6, 1.5) {E};
}
% #1: arestas da árvore, em azul; #2: arestas acrescentadas, em laranja.
\newcommand{\mapaEscolas}[2]{%
\begin{tikzpicture}
  \Mnos
  \foreach \a/\b in {S/A, S/B, A/B, A/D, B/C, B/D, C/E, D/E, C/D}
    \draw [ga] (m\a) -- (m\b);
  \draw [ga] (mS) .. controls (0.5, -1.3) and (3.2, -1.4) .. (mC);
  \foreach \a/\b in {#1} \draw [garv] (m\a) -- (m\b);
  \foreach \a/\b in {#2} \draw [gnova] (m\a) -- (m\b);
  \foreach \a/\b/\p in {S/A/4, S/B/6, A/B/3, A/D/8, B/C/5, B/D/7, C/E/10,
                        D/E/2, C/D/11}
    \path (m\a) -- node [gpeso] {\p} (m\b);
  \path (mS) .. controls (0.5, -1.3) and (3.2, -1.4) .. node [gpeso, pos=0.45] {9} (mC);
  \Mnos
  \node [gv, gpreto] at (mS) {S};
\end{tikzpicture}}

% ------------------------------------------------------------------
% Execução passo a passo de Kruskal e de Prim em G. Os estados seguem a saída
% de kruskal.go e prim.go, e esta parte do arquivo foi gerada por script a
% partir da mesma simulação. Em azul, as arestas da árvore; em laranja, a
% aresta que entra no passo; em vermelho tracejado, a aresta descartada no
% passo.

% Grafo G com pesos no estado de um passo. #1: vértices da árvore (Prim; vazio
% em Kruskal); #2: arestas da árvore; #3: aresta que entra; #4: aresta
% descartada; #5: vértice que acabou de entrar.
\newcommand{\GPestado}[5]{%
  \Gnos
  \Garestas
  \foreach \a/\b in {#2} \draw [garv] (g\a) -- (g\b);
  \foreach \a/\b in {#3} \draw [gnova] (g\a) -- (g\b);
  \foreach \a/\b in {#4} \draw [grej] (g\a) -- (g\b);
  \GPpesos
  \Gnos
  \foreach \v in {#1} \node [gv, gpreto] at (g\v) {\v};
  \foreach \v in {#5} \draw [gatual] (g\v) circle (0.8em);
}

% Um passo de Kruskal. #1 a #4 como em \GPestado (sem #1 e #5); #4: linhas
% da tabela já examinadas, em pares i/situação; #5: linha em exame (0 para
% nenhuma); #6: componentes do Union-Find; #7: o que o passo faz.
\newcommand{\passoKruskal}[7]{%
\begin{tikzpicture}
  \useasboundingbox (-1.1, -2.55) rectangle (8.8, 3.65);
  \GPestado{}{#1}{#2}{#3}{}
  \node [glinha] at (6.55, 3.4) {aresta};
  \node [glegenda] at (7.85, 3.4) {situação};
  \draw [black!60] (6.45, 3.18) -- (8.95, 3.18);
  \foreach \t [count=\i] in {4-5\ \ 1, 1-3\ \ 2, 0-2\ \ 3, 0-1\ \ 4, 2-3\ \ 5,
                             2-4\ \ 6, 0-5\ \ 7, 4-6\ \ 8, 3-6\ \ 9, 6-7\ 12}
    \node [glinha] at (6.55, 3.27 - 0.33*\i) {\t};
  \foreach \i/\s in {#4}
    \node [glegenda] at (7.85, 3.27 - 0.33*\i) {\s};
  \ifnum#5>0
    \draw [gatual] (6.45, 3.27 - 0.33*#5 - 0.16) rectangle (8.95, 3.27 - 0.33*#5 + 0.16);
  \fi
  \node [gtexto] at (-1.1, -0.85) {componentes: \texttt{#6}};
  \node [gtexto] at (-1.1, -1.4) {#7};
  \draw [garv] (-1.0, -2.35) -- (-0.5, -2.35);
  \node [glegenda] at (-0.35, -2.35) {na árvore};
  \draw [gnova] (1.3, -2.35) -- (1.8, -2.35);
  \node [glegenda] at (1.95, -2.35) {entra agora};
  \draw [grej] (3.8, -2.35) -- (4.3, -2.35);
  \node [glegenda] at (4.45, -2.35) {descartada: fecharia um ciclo};
\end{tikzpicture}}

% Um passo de Prim. #1 a #5 como em \GPestado; #6: fila de prioridade em
% ordem de peso, em pares aresta/peso; #7: o que o passo faz.
\newcommand{\passoPrim}[7]{%
\begin{tikzpicture}
  \useasboundingbox (-1.1, -2.55) rectangle (8.8, 3.65);
  \GPestado{#1}{#2}{#3}{#4}{#5}
  \node [glegenda, anchor=center] at (7.9, 3.4) {fila de prioridade};
  \foreach \a/\p [count=\i] in {#6}
    \node [gquadro, fill=white] at (7.9, 3.4 - 0.42*\i) {\a\ (\p)};
  \node [gtexto] at (-1.1, -1.05) {#7};
  \node [gv, gpreto, minimum size=1em] at (-0.9, -2.35) {};
  \node [glegenda] at (-0.65, -2.35) {na árvore};
  \draw [gnova] (1.0, -2.35) -- (1.5, -2.35);
  \node [glegenda] at (1.65, -2.35) {entra agora};
  \draw [grej] (3.4, -2.35) -- (3.9, -2.35);
  \node [glegenda] at (4.05, -2.35) {descartada: os dois extremos na árvore};
\end{tikzpicture}}

\newcommand{\kruskalPasso}[1]{\ifcase#1
\or\passoKruskal{}{}{}{}{0}{\{0\} \{1\} \{2\} \{3\} \{4\} \{5\} \{6\} \{7\}}{início: arestas em ordem de peso; cada vértice é uma componente}
\or\passoKruskal{}{4/5}{}{1/entra}{1}{\{0\} \{1\} \{2\} \{3\} \{4 5\} \{6\} \{7\}}{4-5, peso 1: 4 e 5 em componentes diferentes, entra}
\or\passoKruskal{4/5}{1/3}{}{1/entra, 2/entra}{2}{\{0\} \{1 3\} \{2\} \{4 5\} \{6\} \{7\}}{1-3, peso 2: 1 e 3 em componentes diferentes, entra}
\or\passoKruskal{4/5, 1/3}{0/2}{}{1/entra, 2/entra, 3/entra}{3}{\{0 2\} \{1 3\} \{4 5\} \{6\} \{7\}}{0-2, peso 3: 0 e 2 em componentes diferentes, entra}
\or\passoKruskal{4/5, 1/3, 0/2}{0/1}{}{1/entra, 2/entra, 3/entra, 4/entra}{4}{\{0 1 2 3\} \{4 5\} \{6\} \{7\}}{0-1, peso 4: 0 e 1 em componentes diferentes, entra}
\or\passoKruskal{4/5, 1/3, 0/2, 0/1}{}{2/3}{1/entra, 2/entra, 3/entra, 4/entra, 5/descarta}{5}{\{0 1 2 3\} \{4 5\} \{6\} \{7\}}{2-3, peso 5: 2 e 3 na mesma componente, descarta}
\or\passoKruskal{4/5, 1/3, 0/2, 0/1}{2/4}{}{1/entra, 2/entra, 3/entra, 4/entra, 5/descarta, 6/entra}{6}{\{0 1 2 3 4 5\} \{6\} \{7\}}{2-4, peso 6: 2 e 4 em componentes diferentes, entra}
\or\passoKruskal{4/5, 1/3, 0/2, 0/1, 2/4}{}{0/5}{1/entra, 2/entra, 3/entra, 4/entra, 5/descarta, 6/entra, 7/descarta}{7}{\{0 1 2 3 4 5\} \{6\} \{7\}}{0-5, peso 7: 0 e 5 na mesma componente, descarta}
\or\passoKruskal{4/5, 1/3, 0/2, 0/1, 2/4}{4/6}{}{1/entra, 2/entra, 3/entra, 4/entra, 5/descarta, 6/entra, 7/descarta, 8/entra}{8}{\{0 1 2 3 4 5 6\} \{7\}}{4-6, peso 8: 4 e 6 em componentes diferentes, entra}
\or\passoKruskal{4/5, 1/3, 0/2, 0/1, 2/4, 4/6}{}{3/6}{1/entra, 2/entra, 3/entra, 4/entra, 5/descarta, 6/entra, 7/descarta, 8/entra, 9/descarta}{9}{\{0 1 2 3 4 5 6\} \{7\}}{3-6, peso 9: 3 e 6 na mesma componente, descarta}
\or\passoKruskal{4/5, 1/3, 0/2, 0/1, 2/4, 4/6}{6/7}{}{1/entra, 2/entra, 3/entra, 4/entra, 5/descarta, 6/entra, 7/descarta, 8/entra, 9/descarta, 10/entra}{10}{\{0 1 2 3 4 5 6 7\}}{6-7, peso 12: entra; a árvore tem $V-1 = 7$ arestas, fim}
\fi}
\newcommand{\primPasso}[1]{\ifcase#1
\or\passoPrim{0}{}{}{}{0}{0-2/3, 0-1/4, 0-5/7}{início: 0 entra na árvore, e as arestas de 0 vão para a fila}
\or\passoPrim{0, 2}{}{0/2}{}{2}{0-1/4, 2-3/5, 2-4/6, 0-5/7}{retira 0-2, peso 3: 2 entra, e 2-3 e 2-4 vão para a fila}
\or\passoPrim{0, 2, 1}{0/2}{0/1}{}{1}{1-3/2, 2-3/5, 2-4/6, 0-5/7}{retira 0-1, peso 4: 1 entra, e 1-3 vai para a fila}
\or\passoPrim{0, 2, 1, 3}{0/2, 0/1}{1/3}{}{3}{2-3/5, 2-4/6, 0-5/7, 3-6/9}{retira 1-3, peso 2: 3 entra, e 3-6 vai para a fila}
\or\passoPrim{0, 2, 1, 3}{0/2, 0/1, 1/3}{}{2/3}{}{2-4/6, 0-5/7, 3-6/9}{retira 2-3, peso 5: 2 e 3 já estão na árvore, descarta}
\or\passoPrim{0, 2, 1, 3, 4}{0/2, 0/1, 1/3}{2/4}{}{4}{4-5/1, 0-5/7, 4-6/8, 3-6/9}{retira 2-4, peso 6: 4 entra, e 4-5 e 4-6 vão para a fila}
\or\passoPrim{0, 2, 1, 3, 4, 5}{0/2, 0/1, 1/3, 2/4}{4/5}{}{5}{0-5/7, 4-6/8, 3-6/9}{retira 4-5, peso 1: 5 entra, sem arestas novas para a fila}
\or\passoPrim{0, 2, 1, 3, 4, 5}{0/2, 0/1, 1/3, 2/4, 4/5}{}{0/5}{}{4-6/8, 3-6/9}{retira 0-5, peso 7: 0 e 5 já estão na árvore, descarta}
\or\passoPrim{0, 2, 1, 3, 4, 5, 6}{0/2, 0/1, 1/3, 2/4, 4/5}{4/6}{}{6}{3-6/9, 6-7/12}{retira 4-6, peso 8: 6 entra, e 6-7 vai para a fila}
\or\passoPrim{0, 2, 1, 3, 4, 5, 6}{0/2, 0/1, 1/3, 2/4, 4/5, 4/6}{}{3/6}{}{6-7/12}{retira 3-6, peso 9: 3 e 6 já estão na árvore, descarta}
\or\passoPrim{0, 2, 1, 3, 4, 5, 6, 7}{0/2, 0/1, 1/3, 2/4, 4/5, 4/6}{6/7}{}{7}{}{retira 6-7, peso 12: 7 entra; árvore completa, fim}
\fi}

% Heap de mínimo do exemplo da aula de árvores geradoras mínimas: a árvore
% binária completa e o vetor que a guarda. Compartilhado entre o roteiro
% (figuras.tex) e os slides. Os estados foram conferidos por programa. Usa o
% estilo gv de ../../12/img/grafos.tex, que precisa ser carregado antes.

\tikzset{
  hv/.style = {gv, minimum size=1.7em},
  hmov/.style = {fill=orange!45},
  hcmp/.style = {draw=orange!80!black, line width=1.6pt},
  hsai/.style = {fill=red!15, draw=red!75!black, dashed, line width=1.2pt},
  hcel/.style = {rectangle, draw=black, fill=white, minimum width=0.62cm,
    minimum height=0.55cm, inner sep=0pt, line width=0.6pt,
    font=\sffamily\small},
  hidx/.style = {font=\sffamily\scriptsize, text=black!55, inner sep=1pt},
  htexto/.style = {font=\sffamily\small, text width=7.6cm, align=center,
    execute at begin node={\hyphenpenalty=10000}}
}

% Posições dos índices 0 a 7 na árvore.
\newcommand{\heapPos}{%
  \coordinate (h0) at (0, 0);
  \coordinate (h1) at (-1.8, -0.95);
  \coordinate (h2) at (1.8, -0.95);
  \coordinate (h3) at (-2.7, -1.9);
  \coordinate (h4) at (-0.9, -1.9);
  \coordinate (h5) at (0.9, -1.9);
  \coordinate (h6) at (2.7, -1.9);
  \coordinate (h7) at (-3.2, -2.85);
}

% Um estado do heap. #1: elementos do vetor, do índice 0 em diante; #2:
% índices do elemento que se move, em laranja; #3: índices com que ele é
% comparado, com borda laranja; #4: índice do elemento que sai, em vermelho
% tracejado; #5: legenda.
\newcommand{\heapEstado}[5]{%
  \heapPos
  \foreach \x [count=\i from 0] in {#1} {
    \expandafter\xdef\csname hval\i\endcsname{\x}
    \ifnum\i>0
      \pgfmathtruncatemacro{\p}{(\i - 1) / 2}
      \draw [ga] (h\p) -- (h\i);
    \fi
  }
  \foreach \x [count=\i from 0] in {#1} {
    \node [hv] at (h\i) {\x};
    \node [hidx, anchor=north] at ($(h\i) + (0, -0.31)$) {\i};
    \node [hcel] (c\i) at (-2.17 + 0.62*\i, -3.85) {\x};
    \node [hidx, below=0pt of c\i] {\i};
  }
  \foreach \i in {#2} {
    \node [hv, hmov] at (h\i) {\csname hval\i\endcsname};
    \node [hcel, hmov] at (c\i) {\csname hval\i\endcsname};
  }
  \foreach \i in {#3} {
    \node [hv, hcmp] at (h\i) {\csname hval\i\endcsname};
    \node [hcel, hcmp] at (c\i) {\csname hval\i\endcsname};
  }
  \foreach \i in {#4} {
    \node [hv, hsai] at (h\i) {\csname hval\i\endcsname};
    \node [hcel, hsai] at (c\i) {\csname hval\i\endcsname};
  }
  \node [htexto, anchor=north] at (0, -4.5) {#5};
}

% Passos da inserção do 5 (0 a 3) e da retirada do mínimo (4 a 7).
\newcommand{\heapEstadoN}[1]{\ifcase#1
  \heapEstado{2, 9, 4, 12, 10, 6, 8}{}{}{}%
    {Heap de mínimo com 7 elementos: cada um é menor ou igual aos filhos.}
\or\heapEstado{2, 9, 4, 12, 10, 6, 8, 5}{7}{3}{}%
    {Insere o 5: ele entra no fim do vetor, no índice 7. O pai, no índice 3,
     é 12. Como $5 < 12$, os dois trocam.}
\or\heapEstado{2, 9, 4, 5, 10, 6, 8, 12}{3}{1}{}%
    {O pai do índice 3 é o índice 1, com 9. Como $5 < 9$, os dois trocam.}
\or\heapEstado{2, 5, 4, 9, 10, 6, 8, 12}{1}{0}{}%
    {O pai do índice 1 é a raiz, com 2. Como $2 \le 5$, o 5 para.
     Foram duas trocas.}
\or\heapEstado{2, 5, 4, 9, 10, 6, 8, 12}{7}{}{0}%
    {Retira o menor, o 2, que está na raiz. O último elemento, 12, vai para
     o lugar dele.}
\or\heapEstado{12, 5, 4, 9, 10, 6, 8}{0}{1, 2}{}%
    {Os filhos do 12 são 5 e 4. O menor é o 4, e como $4 < 12$, os dois
     trocam.}
\or\heapEstado{4, 5, 12, 9, 10, 6, 8}{2}{5, 6}{}%
    {Os filhos do 12 agora são 6 e 8. O menor é o 6, e como $6 < 12$, os
     dois trocam.}
\or\heapEstado{4, 5, 6, 9, 10, 12, 8}{5}{}{}%
    {O 12 chegou a uma folha e para. O novo menor, 4, está na raiz.}
\fi}

\newcommand{\heapPasso}[1]{%
\begin{tikzpicture}
  \useasboundingbox (-3.9, 0.45) rectangle (3.9, -5.15);
  \heapEstadoN{#1}
\end{tikzpicture}}

% Quatro passos em grade 2 x 2, para o roteiro. #1: primeiro passo.
\newcommand{\heapQuatro}[1]{%
\begin{tikzpicture}
  \begin{scope}
    \pgfmathtruncatemacro{\k}{#1}\heapEstadoN{\k}
  \end{scope}
  \begin{scope}[xshift=8.4cm]
    \pgfmathtruncatemacro{\k}{#1 + 1}\heapEstadoN{\k}
  \end{scope}
  \begin{scope}[yshift=-6.2cm]
    \pgfmathtruncatemacro{\k}{#1 + 2}\heapEstadoN{\k}
  \end{scope}
  \begin{scope}[xshift=8.4cm, yshift=-6.2cm]
    \pgfmathtruncatemacro{\k}{#1 + 3}\heapEstadoN{\k}
  \end{scope}
\end{tikzpicture}}


\setminted{
  fontsize=\footnotesize,
  breaklines=true,
  breakanywhere=true,
  frame=lines,
  framesep=2mm,
  baselinestretch=1.05,
  tabsize=2
}

\apptocmd{\frame}{}{\justifying}{}

\setbeamertemplate{navigation symbols}{}

% slides de prática e de resposta: fundo e título em cores próprias
\newcommand{\estilopratica}{%
  \setbeamercolor{frametitle}{fg=orange!85!black}%
  \setbeamercolor{background canvas}{bg=orange!8}}
\newcommand{\estiloresposta}{%
  \setbeamercolor{frametitle}{fg=green!45!black}%
  \setbeamercolor{background canvas}{bg=green!6}}

\title{13 - Árvores geradoras mínimas}
\author{Prof. Alex Kutzke}
\date{Novembro de 2026}

\begin{document}

% ==========================================
\begin{frame}
	\begin{center}
		\large{Setor de Educação Profissional e Tecnológica}\\
		Tecnologia em Análise e Desenvolvimento de Sistemas\\
		\vspace{1cm}
		\small{DS143 - Estruturas de Dados II}
	\end{center}
	\titlepage
\end{frame}

% ==========================================
\begin{frame}{Objetivos da aula}
	\footnotesize
	Ao final desta aula você deve ser capaz de:
	\begin{enumerate}
		\item Definir árvore geradora e árvore geradora mínima de um grafo com
		      pesos;
		\item Enunciar a propriedade do corte e usá-la para mostrar que uma
		      aresta pertence à árvore mínima, indicando um corte em que ela é
		      a aresta de travessia mais leve;
		\item Executar à mão os algoritmos de Kruskal e de Prim, registrando a
		      ordem em que as arestas entram e as que são descartadas;
		\item Explicar o papel do Union-Find no Kruskal e o da fila de
		      prioridade no Prim, e o custo de cada algoritmo;
		\item Discutir o que a árvore geradora mínima modela bem num projeto de
		      infraestrutura, e o que ela deixa de fora.
	\end{enumerate}
	\vspace{0.1cm}
	Este conteúdo cai na Prova 2.
\end{frame}

% ==========================================
\begin{frame}{Roteiro}
	\tableofcontents
\end{frame}

% ==========================================
\section{Árvores geradoras}

\begin{frame}[fragile]{Demonstração: fibra óptica para escolas rurais}
	\begin{minted}[fontsize=\scriptsize]{text}
$ go run escolas.go
Tudo ligado à sede                  Árvore geradora mínima
|            K             .G   |   |           .K....       ...G   |
|   A.        .      F   ..     |   |   A..   ..      ...F...       |
|     ...B.   .     . ...       |   |      ..B.                     |
|        ....  .  ....          |   |          ..                   |
|           .... ...            |   |            ..                 |
|             ..S...            |   |              .S...            |
|        ....... .......H       |   |              .    ....H       |
|     C....   .  .    ...       |   |     C.      .          ..     |
|   ....     .    .      ....   |   |   ..  ...  .             ..   |
| D.        E     .          .I |   | D.       .E...             .I |
|                  J            |   |               ...J            |
11 ligações, 129.8 km               11 ligações, 81.6 km
\end{minted}

	\begin{itemize}
		\item município fictício: a sede S já tem conexão, as escolas A a K não;
		\item as duas redes ligam todas as escolas; a da direita usa 63\% do cabo.
	\end{itemize}
\end{frame}

\begin{frame}{Grafo com pesos}
	Cada aresta tem um \textbf{peso}: distância, custo, tempo.

	\begin{center}
		\scalebox{0.85}{\grafoGP}
	\end{center}

	\begin{itemize}
		\item o grafo $G$ da aula de representação e percursos, agora com pesos;
		\item os pesos são números quaisquer, inclusive negativos ou repetidos.
	\end{itemize}
\end{frame}

\begin{frame}[fragile]{Representação em Go}
	\begin{minted}{go}
type Aresta struct {
	V, W int
	Peso float64
}

// Outro devolve o extremo da aresta que não é v.
func (a Aresta) Outro(v int) int {
	if v == a.V {
		return a.W
	}
	return a.V
}

type GrafoP struct {
	adj [][]Aresta // a aresta v-w está na lista de v e na de w
}
\end{minted}
	As listas de adjacência guardam a aresta inteira, e não só o vizinho.
\end{frame}

\begin{frame}{Árvore geradora}
	Uma \textbf{árvore geradora} de um grafo conexo é um subconjunto das arestas
	que liga todos os vértices e não tem ciclo.

	\begin{columns}[c]
		\begin{column}{0.5\textwidth}
			\begin{center}
				\scalebox{0.7}{\grafoGPbfs}
			\end{center}
		\end{column}
		\begin{column}{0.5\textwidth}
			\small
			\begin{itemize}
				\item tem sempre $V - 1$ arestas;
				\item a árvore da busca em largura é uma árvore geradora;
				\item à esquerda, a de $G$ a partir de 0, com peso 43:
				      \\ $4 + 3 + 7 + 2 + 6 + 9 + 12$.
			\end{itemize}
		\end{column}
	\end{columns}
\end{frame}

\begin{frame}{Árvore geradora mínima}
	A \textbf{árvore geradora mínima} (AGM, \emph{minimum spanning tree}) é a
	árvore geradora de menor peso total.

	\begin{columns}[c]
		\begin{column}{0.5\textwidth}
			\begin{center}
				\scalebox{0.7}{\grafoGPagm}
			\end{center}
		\end{column}
		\begin{column}{0.5\textwidth}
			\small
			\begin{itemize}
				\item em $G$, peso 36;
				\item $G$ tem 50 árvores geradoras; o grafo completo com 12
				      vértices da demonstração tem $12^{10}$, cerca de 62 bilhões;
				\item testar todas está fora de questão.
			\end{itemize}
		\end{column}
	\end{columns}

	\vspace{0.2cm}
	\footnotesize
	Com pesos distintos, a árvore mínima é única. Com pesos repetidos, pode
	haver mais de uma, todas com o mesmo peso.
\end{frame}

\begin{frame}{Propriedade do corte}
	\small
	Um \textbf{corte} divide os vértices em dois conjuntos não vazios. Uma
	\textbf{aresta de travessia} liga um vértice de cada lado.

	\begin{columns}[c]
		\begin{column}{0.55\textwidth}
			\begin{center}
				\scalebox{0.65}{\grafoGPcorte}
			\end{center}
		\end{column}
		\begin{column}{0.45\textwidth}
			\footnotesize
			Arestas de travessia: 0-5 (7), 2-4 (6) e 3-6 (9).

			\vspace{0.2cm}
			\textbf{Propriedade do corte}: em qualquer corte, a aresta de
			travessia de menor peso está na árvore geradora mínima (pesos
			distintos).
		\end{column}
	\end{columns}
\end{frame}

\begin{frame}{Por que a propriedade do corte vale}
	\begin{columns}[T,onlytextwidth]
		\begin{column}{0.44\textwidth}
			\footnotesize
			Seja $e$ a aresta de travessia mais leve, e suponha uma árvore
			mínima $T$ sem $e$.
			\begin{enumerate}
				\item acrescentar $e$ a $T$ fecha um ciclo;
				\item o ciclo atravessa o corte por $e$ e precisa voltar: alguma
				      outra aresta $f$ do ciclo também é de travessia;
				\item trocar $f$ por $e$ dá outra árvore geradora;
				\item $e$ pesa menos que $f$, então a árvore nova pesa menos que
				      $T$.
			\end{enumerate}
			Contradição: $T$ não era mínima.
		\end{column}
		\begin{column}{0.54\textwidth}
			\centering
			\scalebox{0.83}{%
				\only<1>{\provaCorte{1}}%
				\only<2>{\provaCorte{2}}%
				\only<3>{\provaCorte{3}}}
		\end{column}
	\end{columns}
\end{frame}

\begin{frame}{Algoritmo guloso}
	\begin{enumerate}
		\item comece sem nenhuma aresta;
		\item escolha um corte que nenhuma aresta já escolhida atravessa;
		\item acrescente a aresta de travessia mais leve desse corte;
		\item repita até ter $V - 1$ arestas.
	\end{enumerate}

	\vspace{0.3cm}
	Cada aresta acrescentada está na árvore mínima, pela propriedade do corte.
	Kruskal e Prim são duas formas de escolher o corte a cada passo.
\end{frame}

% ==========================================
\section{Kruskal}

\begin{frame}{Algoritmo de Kruskal}
	\begin{enumerate}
		\item ordene as arestas por peso, da mais leve para a mais pesada;
		\item examine cada aresta $v$-$w$ nessa ordem:
		      \begin{itemize}
			      \item se $v$ e $w$ já estão ligados pelas arestas escolhidas, a
			            aresta fecharia um ciclo: descarte;
			      \item se não, a aresta entra;
		      \end{itemize}
		\item pare com $V - 1$ arestas.
	\end{enumerate}

	\vspace{0.3cm}
	``$v$ e $w$ já estão ligados?'' é a pergunta do Union-Find, a primeira
	estrutura do semestre: \texttt{Connected} e \texttt{Union}, na versão com
	união por tamanho.
\end{frame}

\begin{frame}[fragile]{Kruskal em Go}
	\begin{minted}[fontsize=\scriptsize]{go}
func Kruskal(n int, arestas []Aresta) []Aresta {
	ordem := make([]Aresta, len(arestas))
	copy(ordem, arestas)
	sort.Slice(ordem, func(i, j int) bool {
		return ordem[i].Peso < ordem[j].Peso
	})

	uf := NovoUF(n)
	var arvore []Aresta
	for _, a := range ordem {
		if len(arvore) == n-1 {
			break // a árvore já tem V-1 arestas
		}
		if uf.Connected(a.V, a.W) {
			continue // fecharia um ciclo
		}
		uf.Union(a.V, a.W)
		arvore = append(arvore, a)
	}
	return arvore
}
\end{minted}
\end{frame}

\begin{frame}{Kruskal passo a passo}
	\begin{center}
		\only<1>{\kruskalPasso{1}}
		\only<2>{\kruskalPasso{2}}
		\only<3>{\kruskalPasso{3}}
		\only<4>{\kruskalPasso{4}}
		\only<5>{\kruskalPasso{5}}
		\only<6>{\kruskalPasso{6}}
		\only<7>{\kruskalPasso{7}}
		\only<8>{\kruskalPasso{8}}
		\only<9>{\kruskalPasso{9}}
		\only<10>{\kruskalPasso{10}}
		\only<11>{\kruskalPasso{11}}
	\end{center}
\end{frame}

\begin{frame}[fragile]{Kruskal em $G$}
	\begin{minted}[fontsize=\scriptsize]{text}
$ go run kruskal.go
Kruskal: arestas da mais leve para a mais pesada
  4-5     1   entra
  1-3     2   entra
  0-2     3   entra
  0-1     4   entra
  2-3     5   descarta: 2 e 3 já estão ligados
  2-4     6   entra
  0-5     7   descarta: 0 e 5 já estão ligados
  4-6     8   entra
  3-6     9   descarta: 3 e 6 já estão ligados
  6-7    12   entra

Árvore geradora mínima: 4-5 1-3 0-2 0-1 2-4 4-6 6-7
7 arestas, peso total 36
\end{minted}
	\footnotesize
	Custo: ordenar as arestas, $O(E \log E)$, mais até $E$ consultas ao
	Union-Find, $O(\log V)$ cada. Total $O(E \log E)$.
\end{frame}

% ==========================================
\section{Prim}

\begin{frame}{Algoritmo de Prim}
	\begin{enumerate}
		\item comece a árvore por um vértice $s$;
		\item acrescente a aresta mais leve que liga a árvore a um vértice fora
		      dela, e esse vértice;
		\item repita até a árvore ter todos os vértices.
	\end{enumerate}

	\vspace{0.3cm}
	O corte de cada passo é \{vértices da árvore\} e \{os demais\}.

	\vspace{0.3cm}
	Para achar a aresta mais leve sem percorrer todas a cada passo, as
	candidatas ficam numa \textbf{fila de prioridade}.
\end{frame}

\begin{frame}{Heap de mínimo}
	\begin{columns}[c,onlytextwidth]
		\begin{column}{0.56\textwidth}
			\centering
			\scalebox{0.74}{\heapPasso{0}}
		\end{column}
		\begin{column}{0.42\textwidth}
			\footnotesize
			Fila de prioridade: insere em qualquer ordem e retira sempre o
			menor.
			\begin{itemize}
				\item árvore binária completa, com os níveis preenchidos da
				      esquerda para a direita;
				\item cada elemento é menor ou igual aos filhos, então o menor
				      está na raiz;
				\item guardada num vetor: os filhos de $i$ estão em $2i+1$ e
				      $2i+2$, e o pai em $(i-1)/2$.
			\end{itemize}
			A altura é $\lfloor \log_2 n \rfloor$: inserir e retirar o menor
			custam $O(\log n)$.
		\end{column}
	\end{columns}
\end{frame}

\begin{frame}{Inserção: o novo elemento sobe}
	\begin{center}
		\only<1>{\heapPasso{0}}
		\only<2>{\heapPasso{1}}
		\only<3>{\heapPasso{2}}
		\only<4>{\heapPasso{3}}
	\end{center}
\end{frame}

\begin{frame}[fragile]{Inserção em Go}
	\begin{columns}[T,onlytextwidth]
		\begin{column}{0.44\textwidth}
			\begin{minted}[fontsize=\fontsize{7}{8.5}\selectfont]{go}
type HeapMin struct {
	v []int
}

func (h *HeapMin) Insere(x int) {
	h.v = append(h.v, x)
	h.sobe(len(h.v) - 1)
}
\end{minted}
		\end{column}
		\begin{column}{0.54\textwidth}
			\begin{minted}[fontsize=\fontsize{7}{8.5}\selectfont]{go}
func (h *HeapMin) sobe(i int) {
	for i > 0 {
		pai := (i - 1) / 2
		if h.v[pai] <= h.v[i] {
			return
		}
		h.v[pai], h.v[i] = h.v[i], h.v[pai]
		i = pai
	}
}
\end{minted}
		\end{column}
	\end{columns}
	\footnotesize
	O novo elemento entra no fim do vetor e troca de lugar com o pai enquanto
	for menor que ele. São no máximo $\lfloor \log_2 n \rfloor$ trocas, uma
	por nível.
\end{frame}

\begin{frame}{Retirada do mínimo: o último elemento desce}
	\begin{center}
		\only<1>{\heapPasso{4}}
		\only<2>{\heapPasso{5}}
		\only<3>{\heapPasso{6}}
		\only<4>{\heapPasso{7}}
	\end{center}
\end{frame}

\begin{frame}[fragile]{Retirada do mínimo em Go}
	\begin{columns}[T,onlytextwidth]
		\begin{column}{0.44\textwidth}
			\begin{minted}[fontsize=\fontsize{7}{8.5}\selectfont]{go}
func (h *HeapMin) RetiraMin() int {
	menor := h.v[0]
	ultimo := len(h.v) - 1
	h.v[0] = h.v[ultimo]
	h.v = h.v[:ultimo]
	if len(h.v) > 0 {
		h.desce(0)
	}
	return menor
}
\end{minted}
		\end{column}
		\begin{column}{0.54\textwidth}
			\begin{minted}[fontsize=\fontsize{7}{8.5}\selectfont]{go}
func (h *HeapMin) desce(i int) {
	n := len(h.v)
	for {
		menor := i
		esq, dir := 2*i+1, 2*i+2
		if esq < n && h.v[esq] < h.v[menor] {
			menor = esq
		}
		if dir < n && h.v[dir] < h.v[menor] {
			menor = dir
		}
		if menor == i {
			return
		}
		h.v[menor], h.v[i] = h.v[i], h.v[menor]
		i = menor
	}
}
\end{minted}
		\end{column}
	\end{columns}
	\footnotesize
	O último elemento vai para a raiz e troca de lugar com o menor dos filhos
	enquanto for maior que ele. Também são no máximo
	$\lfloor \log_2 n \rfloor$ trocas.
\end{frame}

\begin{frame}[fragile]{Fila de prioridade com \texttt{container/heap}}
	\small
	O pacote faz o sobe e o desce. O tipo informa o tamanho, a comparação e a
	troca, e como acrescentar e retirar do fim do \textit{slice}.

	\begin{minted}[fontsize=\fontsize{7}{8.5}\selectfont]{go}
// O pacote container/heap organiza o slice como heap de mínimo.
type FilaArestas []Aresta

func (f FilaArestas) Len() int { return len(f) }
func (f FilaArestas) Less(i, j int) bool { return f[i].Peso < f[j].Peso }
func (f FilaArestas) Swap(i, j int) { f[i], f[j] = f[j], f[i] }
func (f *FilaArestas) Push(x any) { *f = append(*f, x.(Aresta)) }
func (f *FilaArestas) Pop() any {
	antigo := *f
	n := len(antigo)
	a := antigo[n-1]
	*f = antigo[:n-1]
	return a
}
\end{minted}
	\footnotesize
	O programa usa \texttt{heap.Push(fila, a)} e \texttt{heap.Pop(fila)}, e não
	os métodos diretamente. No trabalho prático, \texttt{container/*} é
	proibido.
\end{frame}

\begin{frame}[fragile]{Prim em Go}
	\begin{columns}[T,onlytextwidth]
		\begin{column}{0.5\textwidth}
			\begin{minted}[fontsize=\fontsize{7}{8.5}\selectfont]{go}
func Prim(g *GrafoP, s int) []Aresta {
	naArvore := make([]bool, g.V())
	fila := &FilaArestas{}
	var arvore []Aresta

	// v entra na árvore, e as
	// arestas de v para fora dela,
	// na fila.
	visita := func(v int) {
		naArvore[v] = true
		for _, a := range g.adj[v] {
			if !naArvore[a.Outro(v)] {
				heap.Push(fila, a)
			}
		}
	}
\end{minted}
		\end{column}
		\begin{column}{0.5\textwidth}
			\begin{minted}[fontsize=\fontsize{7}{8.5}\selectfont]{go}
	visita(s)
	for fila.Len() > 0 &&
		len(arvore) < g.V()-1 {
		a := heap.Pop(fila).(Aresta)
		if naArvore[a.V] && naArvore[a.W] {
			continue // descarta
		}
		arvore = append(arvore, a)
		novo := a.V
		if naArvore[a.V] {
			novo = a.W
		}
		visita(novo)
	}
	return arvore
}
\end{minted}
		\end{column}
	\end{columns}
\end{frame}

\begin{frame}{Prim passo a passo}
	\begin{center}
		\only<1>{\primPasso{1}}
		\only<2>{\primPasso{2}}
		\only<3>{\primPasso{3}}
		\only<4>{\primPasso{4}}
		\only<5>{\primPasso{5}}
		\only<6>{\primPasso{6}}
		\only<7>{\primPasso{7}}
		\only<8>{\primPasso{8}}
		\only<9>{\primPasso{9}}
		\only<10>{\primPasso{10}}
		\only<11>{\primPasso{11}}
	\end{center}
\end{frame}

\begin{frame}[fragile]{Prim em $G$, a partir de 0}
	\begin{minted}[fontsize=\scriptsize]{text}
$ go run prim.go
Prim a partir de 0
  início            fila: 0-2(3) 0-1(4) 0-5(7)
  0-2  3 entra 2    fila: 0-1(4) 2-3(5) 2-4(6) 0-5(7)
  0-1  4 entra 1    fila: 1-3(2) 2-3(5) 2-4(6) 0-5(7)
  1-3  2 entra 3    fila: 2-3(5) 2-4(6) 0-5(7) 3-6(9)
  2-3  5 descarta   fila: 2-4(6) 0-5(7) 3-6(9)
  2-4  6 entra 4    fila: 4-5(1) 0-5(7) 4-6(8) 3-6(9)
  4-5  1 entra 5    fila: 0-5(7) 4-6(8) 3-6(9)
  0-5  7 descarta   fila: 4-6(8) 3-6(9)
  4-6  8 entra 6    fila: 3-6(9) 6-7(12)
  3-6  9 descarta   fila: 6-7(12)
  6-7 12 entra 7    fila:

Árvore geradora mínima: 0-2 0-1 1-3 2-4 4-5 4-6 6-7
7 arestas, peso total 36
\end{minted}
	\footnotesize
	A mesma árvore do Kruskal, em outra ordem. A aresta 4-5, a mais leve, é
	a quinta a entrar.
\end{frame}

\begin{frame}{Custo do Prim}
	\begin{itemize}
		\item cada aresta entra na fila no máximo uma vez e sai no máximo uma
		      vez: $O(E \log E)$, com até $E$ arestas na fila;
		\item esta é a versão \textbf{preguiçosa}: arestas que deixaram de servir
		      ficam na fila e são descartadas quando saem;
		\item a versão \textbf{ansiosa} guarda na fila um vértice por vez, com a
		      melhor aresta conhecida até ele: $O(E \log V)$, com no máximo $V$
		      itens na fila;
		\item em grafo denso, um vetor no lugar da fila dá $O(V^2)$, sem heap.
	\end{itemize}
\end{frame}

\begin{frame}{Kruskal ou Prim}
	\small
	\begin{center}
		\begin{tabular}{>{\raggedright}p{2.6cm}>{\raggedright}p{3.4cm}>{\raggedright\arraybackslash}p{3.4cm}}
			\toprule
			& Kruskal & Prim \\
			\midrule
			o que cresce & uma floresta, que vira árvore no fim & uma árvore, a partir de $s$ \\
			estrutura auxiliar & ordenação e Union-Find & fila de prioridade \\
			custo & $O(E \log E)$ & $O(E \log E)$, ou $O(E \log V)$ na versão ansiosa \\
			grafo desconexo & floresta geradora mínima & só a componente de $s$ \\
			\bottomrule
		\end{tabular}
	\end{center}

	\vspace{0.3cm}
	Como $E < V^2$, $\log E < 2 \log V$: $O(E \log E)$ e $O(E \log V)$ diferem
	por uma constante.
\end{frame}

% ==========================================
\section{Infraestrutura}

\begin{frame}{Origem do problema}
	\small
	\begin{itemize}
		\item 1926, Otakar Bor\r{u}vka: rede elétrica da Morávia, na atual República
		      Tcheca, com o menor comprimento de linhas;
		\item 1930, Vojt\v{e}ch Jarník: o algoritmo que hoje se chama de Prim;
		\item 1956, Joseph Kruskal; 1957, Robert Prim, nos Bell Labs, para redes
		      telefônicas;
		\item 1959, Edsger Dijkstra redescobre o algoritmo de Jarník e Prim.
	\end{itemize}

	\vspace{0.2cm}
	ODS 9, Indústria, Inovação e Infraestrutura:
	\begin{itemize}
		\footnotesize
		\item[9.1] infraestrutura de qualidade, confiável, sustentável e
		      \textbf{resiliente}, com acesso equitativo e a preços acessíveis;
		\item[9.c] aumentar o acesso às tecnologias de informação e comunicação
		      e oferecer acesso universal à internet.
	\end{itemize}
\end{frame}

{\estilopratica
\begin{frame}{Prática: a rede das escolas}
	\small
	Custo de cada ligação possível, em km de cabo. A sede S já tem conexão.

	\begin{center}
		\scalebox{0.8}{\mapaEscolas{}{}}
	\end{center}

	\begin{enumerate}
		\item execute o Kruskal à mão: liste as arestas na ordem em que são
		      examinadas e marque as que entram e as que são descartadas;
		\item calcule o custo total da rede.
	\end{enumerate}
	Em dupla. Tempo: 7 minutos.
\end{frame}
}

{\estiloresposta
\begin{frame}{Resposta: a rede das escolas}
	\small
	\begin{columns}[c]
		\begin{column}{0.55\textwidth}
			\begin{center}
				\scalebox{0.7}{\mapaEscolas{S/A, A/B, B/C, B/D, D/E}{}}
			\end{center}
		\end{column}
		\begin{column}{0.45\textwidth}
			\footnotesize
			D-E 2 entra\\
			A-B 3 entra\\
			S-A 4 entra\\
			B-C 5 entra\\
			S-B 6 descarta\\
			B-D 7 entra: 5 arestas, fim\\[0.2cm]
			Custo: $2 + 3 + 4 + 5 + 7 = 21$ km.
		\end{column}
	\end{columns}

	\vspace{0.2cm}
	\footnotesize
	A-D, S-C, C-E e C-D nem chegam a ser examinadas.
\end{frame}
}

{\estilopratica
\begin{frame}{Prática: e se um cabo for cortado?}
	\small
	Na rede de 21 km:
	\begin{enumerate}
		\item qual ligação, se for cortada, deixa mais escolas sem conexão com a
		      sede? Quantas?
		\item acrescente uma única ligação, a mais barata possível, de modo que
		      nenhum corte isole mais de duas escolas;
		\item quais ligações é preciso acrescentar para que nenhum corte isole
		      escola alguma? Qual o custo total da rede?
	\end{enumerate}
	Em dupla. Tempo: 8 minutos.
\end{frame}
}

{\estiloresposta
\begin{frame}{Resposta: e se um cabo for cortado?}
	\footnotesize
	\begin{columns}[c]
		\begin{column}{0.5\textwidth}
			\begin{center}
				\scalebox{0.6}{\mapaEscolas{S/A, A/B, B/C, B/D, D/E}{S/B, C/E}}
			\end{center}
		\end{column}
		\begin{column}{0.5\textwidth}
			\begin{enumerate}
				\item S-A: as cinco escolas;
				\item S-B, 6 km. O pior corte passa a ser B-D, que isola D e E;
				\item S-B e C-E: 16 km a mais, rede de 37 km, 76\% mais cara.
			\end{enumerate}
		\end{column}
	\end{columns}

	\vspace{0.2cm}
	Na árvore, toda aresta é uma ponte: retirada, desconecta o grafo. A
	redundância custa caro. A rede mais barata sem ponte custa 35 km, e não
	contém A-B nem B-D: partir da árvore e acrescentar arestas não dá o ótimo.
\end{frame}
}

\begin{frame}{O que o modelo deixa de fora}
	\small
	\begin{itemize}
		\item \textbf{resiliência}: a árvore mínima não tem nenhuma redundância,
		      e a meta 9.1 pede infraestrutura resiliente;
		\item \textbf{pontos de junção}: com três escolas nos vértices de um
		      triângulo de 10 km de lado, a árvore mínima usa 20 km; uma caixa de
		      emenda no centro liga as três com 17,3 km. Escolher os pontos
		      extras é o problema da árvore de Steiner, NP-difícil;
		\item \textbf{custo real}: travessia de rio, licença ambiental, poste já
		      existente. O peso da aresta é uma decisão de modelagem;
		\item \textbf{quem fica de fora}: a árvore liga todos os vértices que
		      estão no grafo. Quem decide quais escolas entram é o projeto, e não
		      o algoritmo.
	\end{itemize}
\end{frame}

\begin{frame}{Onde aparecem as árvores geradoras mínimas}
	\footnotesize
	\begin{tabular}{>{\raggedright}p{3.6cm}>{\raggedright\arraybackslash}p{6.4cm}}
		\toprule
		Aplicação & O que a árvore representa \\
		\midrule
		projeto de redes & elétrica, telefônica, de dados, de água, viária \\
		agrupamento (\emph{single-link}) & tirar as $k-1$ arestas mais pesadas
		da árvore deixa $k$ grupos \\
		segmentação de imagens & pixels como vértices, diferença de cor como
		peso \\
		caixeiro-viajante & a árvore dá uma rota no máximo duas vezes maior que a
		ótima, com pesos que obedecem à desigualdade triangular \\
		taxonomia e biologia & árvores de semelhança entre espécies ou genes \\
		\bottomrule
	\end{tabular}

	\vspace{0.2cm}
	\scriptsize
	Lista baseada nos slides de Sedgewick e Wayne, \emph{Minimum Spanning
	Trees}.
\end{frame}

% ==========================================
\section{Fechamento}

\begin{frame}[fragile]{De volta às escolas}
	\begin{minted}[fontsize=\scriptsize]{text}
Ligações da árvore e escolas isoladas da sede se a ligação for cortada:
  A-B   5.4 km   A
  B-K   5.7 km   F G K
  C-D   5.7 km   D
  C-E   7.2 km   C D
  E-J   7.3 km   J
  F-G   7.3 km   G
  F-K   8.2 km   F G
  S-H   8.2 km   H I
  H-I   8.5 km   I
  S-E   8.9 km   C D E J
  S-B   9.2 km   A B F G K
\end{minted}
	\footnotesize
	\begin{itemize}
		\item 66 ligações possíveis, Kruskal sobre o grafo completo;
		\item o corte de S-B deixa cinco escolas sem conexão.
	\end{itemize}
\end{frame}

\begin{frame}{Para a próxima aula}
	\begin{itemize}
		\item exercícios do roteiro da aula, sem entrega e sem nota. O conteúdo
		      cai na Prova 2;
		\item rodar \texttt{kruskal.go}, \texttt{prim.go} e
		      \texttt{escolas.go}, trocando as arestas e os locais;
		\item leitura: Sedgewick, seção 4.3; Cormen, capítulo 23.
	\end{itemize}

	\vspace{0.4cm}
	Próxima aula, 11/11: caminhos mínimos, com a Atividade Avaliativa 4. O
	algoritmo de Dijkstra tem a estrutura do Prim, com outra prioridade na fila.
	No mesmo dia vence a etapa 1 do trabalho prático.
\end{frame}

\end{document}
